\section{Motivation}

A matrix acting on a space moves every vector somewhere else, and at first sight the motion is a mess.
Eigenvectors are the directions in which the mess disappears: along an eigenvector the matrix acts as mere
scaling by a number — the eigenvalue. Finding these special directions means finding the natural
coordinate system of a transformation, the axes along which its behaviour is simplest. It is hard to
overstate how often computing profits from this one idea:

\begin{itemize}
\item \textbf{Ranking the web.} Google's founding insight was that the importance of a web page is an
eigenvector computation: importance flows along links, and a stable assignment of importances is precisely
an eigenvector of the web's matrix for the eigenvalue $1$. The final section of this chapter works this
out in detail.
\item \textbf{Dimensionality reduction.} Principal Component Analysis (PCA), a staple of data science,
computes the eigenvectors of a data set's covariance matrix: they are the directions along which the data
genuinely varies, and the eigenvalues say by how much. Keeping only the few directions with the largest
eigenvalues compresses the data with minimal loss — the classic ``eigenfaces'' method recognised human
faces this way.
\item \textbf{Repeated processes.} Whatever evolves in discrete steps by a fixed linear rule — a Markov
chain, a population model, the pair of previous terms in the Fibonacci recurrence — evolves by matrix
powers $A^k$. Eigenvalues make such powers computable in closed form (see the section on diagonalisation),
and the largest eigenvalue governs the long-run fate of the process: growth, decay or convergence to a
steady state. In essence this is also why eigenvalues decide the stability of physical systems, and why
they appear when engineers ask whether a bridge will resonate in the wind.
\item \textbf{Graphs and clustering.} The eigenvalues of matrices associated with a graph reveal its
global shape: how well connected it is and whether it splits into communities. \emph{Spectral clustering}
and much of modern network analysis are built on this.
\end{itemize}

Within this book, eigenvalues are also a grand finale in a literal sense: computing them exercises
everything that came before. The eigenvalue condition is a determinant equation
(Chapter~\ref{ch:determinants}); its solutions are the roots of a polynomial, which may force us into the
complex numbers of Chapter~\ref{ch:complex}; and finding the eigenvectors for each eigenvalue means
solving a homogeneous linear system (Chapter~\ref{ch:sole}) whose solution set is a subspace
(Chapter~\ref{ch:vector-spaces}).

\section{Eigenvalues and eigenvectors}

\begin{definition}
Let $A$ be an $n\times n$ matrix over the field $\K$. A scalar $\lambda\in \K$ is called an \emph{eigenvalue}
of $A$ if there is a non-zero vector $\vv\in \K^n$ such that
\[
A\vv= \lambda \cdot \vv.
\]
\end{definition}

In computations one uses the following equivalent characterisation.

\begin{theorem}
Let $A$ be an $n\times n$ matrix over the field $\K$. A scalar $\lambda\in \K$ is an eigenvalue of $A$ if
and only if
\[
\det(A-\lambda \cdot I) = 0.
\]
\end{theorem}

\begin{proof}
($\Rightarrow$) Let $\vv$ be a non-zero vector with $A\vv = \lambda\cdot \vv$. Then
\[
A\vv - \lambda\cdot \vv = \0, \qquad (A - \lambda \cdot I)\cdot \vv = \0,
\]
so the system $(A - \lambda \cdot I)X = \0$ has a non-zero solution. This happens precisely when
$\det(A - \lambda \cdot I)=0$.

($\Leftarrow$) Assume $\det(A - \lambda \cdot I)=0$. Then the system $(A - \lambda \cdot I)X = \0$ has a
non-zero solution $\vv$, and
\[
(A - \lambda \cdot I)\cdot \vv = \0, \qquad A\vv - \lambda\cdot \vv = \0, \qquad A\vv = \lambda\cdot \vv. \qedhere
\]
\end{proof}

\begin{definition}
The polynomial
\[
p_A(\lambda) = \det(A-\lambda\cdot I)
\]
(a polynomial of degree $n$ in the variable $\lambda$) is called the \emph{characteristic polynomial} of
$A$. The eigenvalues of $A$ are exactly the roots of $p_A$ lying in $\K$.
\end{definition}

\begin{fact}
Similar matrices have the same characteristic polynomial: if $B = P^{-1}AP$, then
\[
p_B(\lambda) = \det(P^{-1}AP - \lambda I) = \det\big(P^{-1}(A-\lambda I)P\big)
= \det(P^{-1})\cdot\det(A-\lambda I)\cdot\det(P) = p_A(\lambda).
\]
\end{fact}

\begin{remark}[Eigenvalues of a linear operator]
For a linear operator $F\colon\V\to\V$, eigenvalues and eigenvectors are defined by the same condition
$F(\vv) = \lambda\cdot\vv$ (with $\vv\neq\0$). They are computed from the matrix of $F$ in \emph{any}
basis $S$: the matrices $M_S(F)$ in different bases are similar (Chapter~\ref{ch:linear-mappings}), so by
the fact above they all have the same characteristic polynomial and hence the same eigenvalues.
\end{remark}

\begin{definition}
A non-zero vector $\vv\in \K^n$ is called an \emph{eigenvector} for the eigenvalue $\lambda$ if
$A\cdot \vv = \lambda \cdot \vv$.
\end{definition}

\begin{theorem}
Let $\lambda$ be an eigenvalue of $A$. The set $W_\lambda$ of all eigenvectors for $\lambda$ (together with
$\0$) is a subspace of $\K^n$.
\end{theorem}

\begin{proof}
Note that $\0\in W_\lambda$. For $\vv_1,\vv_2\in W_\lambda$,
\[
A(\vv_1+\vv_2)=A\vv_1+A\vv_2=\lambda\cdot \vv_1+\lambda\cdot \vv_2 = \lambda \cdot (\vv_1+\vv_2),
\]
so $\vv_1+\vv_2\in W_\lambda$. Similarly, for $\vv\in W_\lambda$ and $k\in \K$ we have
$A(k\vv)=kA\vv=k\lambda\vv=\lambda(k\vv)$, so $k\cdot \vv\in W_\lambda$.
\end{proof}

\section{Computing eigenvalues and eigenvectors}

The recipe is always the same: find the roots of the characteristic polynomial $p_A(\lambda)$, and for each
root $\lambda$ solve the homogeneous system $(A-\lambda I)X = \0$ to find the eigenvectors.

\begin{example}[A $2\times 2$ matrix]
Let $A = \left( \begin{smallmatrix} 2 & 1\\ 1 & 2 \end{smallmatrix} \right)$ over $\R$. Then
\[
p_A(\lambda) = \left| \begin{array}{cc} 2-\lambda & 1\\ 1 & 2-\lambda \end{array} \right|
= (2-\lambda)^2 - 1 = \lambda^2-4\lambda+3 = (\lambda-1)(\lambda-3),
\]
so the eigenvalues are $\lambda = 1$ and $\lambda = 3$.

\medskip
\noindent\emph{Eigenvectors for $\lambda = 1$.} The system $(A-I)X = \0$ reads $x+y = 0$ (twice), so
\[
W_1 = \Big\{ \left( \begin{array}{c} t \\ -t \end{array} \right) \;\Big|\; t\in \R \Big\}.
\]

\noindent\emph{Eigenvectors for $\lambda = 3$.} The system $(A-3I)X = \0$ reads $-x+y = 0$, so
\[
W_3 = \Big\{ \left( \begin{array}{c} t \\ t \end{array} \right) \;\Big|\; t\in \R \Big\}.
\]
Indeed, $A\left( \begin{smallmatrix} 1\\ 1 \end{smallmatrix} \right)
= \left( \begin{smallmatrix} 3\\ 3 \end{smallmatrix} \right)
= 3 \left( \begin{smallmatrix} 1\\ 1 \end{smallmatrix} \right)$.
\end{example}

\begin{example}[A $3\times 3$ matrix]
Let
\[
A=\left( \begin{array}{ccc} 1 & 0 & 2\\ 0 & -1 & 3\\ 0 & 1 & 1 \end{array} \right),
\qquad
A-\lambda \cdot I=\left( \begin{array}{ccc} 1-\lambda & 0 & 2\\ 0 & -1-\lambda & 3\\ 0 & 1 & 1-\lambda \end{array} \right).
\]
Expanding along the first column,
\[
\det(A-\lambda \cdot I)=
(1-\lambda)\cdot \left| \begin{array}{cc} -1-\lambda & 3\\ 1 & 1-\lambda \end{array} \right|
= (1-\lambda)\big((-1-\lambda)(1-\lambda)-3\big)
= (1-\lambda)(\lambda^2-4) = (1-\lambda)(\lambda-2)(\lambda+2).
\]
The eigenvalues of $A$ are therefore $\lambda = 1,\ 2,\ -2$.

\medskip
\noindent\emph{Eigenvectors for $\lambda = 1$.} They are the solutions of
\[
\left( \begin{array}{ccc} 0 & 0 & 2\\ 0 & -2 & 3\\ 0 & 1 & 0 \end{array} \right)X = \0
\xrightarrow{\ \text{row reduction}\ }
\left( \begin{array}{ccc} 0 & 1 & 0 \\ 0 & 0 & 1\\ 0 & 0 & 0 \end{array} \right)X = \0,
\qquad
W_1 = \Big\{ \left( \begin{array}{c} x \\ 0 \\ 0 \end{array} \right) \;\Big|\; x\in \R \Big\}.
\]

\noindent\emph{Eigenvectors for $\lambda = 2$.} They are the solutions of
\[
\left( \begin{array}{ccc} -1 & 0 & 2\\ 0 & -3 & 3\\ 0 & 1 & -1 \end{array} \right)X = \0
\xrightarrow{\ \text{row reduction}\ }
\left( \begin{array}{ccc} 1 & 0 & -2\\ 0 & 1 & -1\\ 0 & 0 & 0 \end{array} \right)X = \0,
\qquad
W_2 = \Big\{ \left( \begin{array}{c} 2z \\ z \\ z \end{array} \right) \;\Big|\; z\in \R \Big\}.
\]

\noindent\emph{Eigenvectors for $\lambda = -2$.} They are the solutions of
\[
\left( \begin{array}{ccc} 3 & 0 & 2\\ 0 & 1 & 3\\ 0 & 1 & 3 \end{array} \right)X = \0
\xrightarrow{\ \text{row reduction}\ }
\left( \begin{array}{ccc} 1 & 0 & \tfrac{2}{3}\\ 0 & 1 & 3\\ 0 & 0 & 0 \end{array} \right)X = \0,
\qquad
W_{-2} = \Big\{ \left( \begin{array}{c} -\tfrac{2}{3}z \\ -3z \\ z \end{array} \right) \;\Big|\; z\in \R \Big\}.
\]
\end{example}

\begin{fact}
For an $n\times n$ matrix $A$ there are at most $n$ distinct eigenvalues of $A$ (the characteristic
polynomial has degree $n$, so it has at most $n$ roots).
\end{fact}

\begin{fact}
The eigenvalues of an upper (or lower) triangular matrix are exactly its diagonal entries: for a triangular
$A$ the matrix $A-\lambda I$ is again triangular, so
\[
p_A(\lambda) = (a_{11}-\lambda)(a_{22}-\lambda)\cdots(a_{nn}-\lambda).
\]
\end{fact}

\begin{remark}
It may happen that a matrix $A$ has no eigenvalues over a field $\K$, or that some eigenvalues are multiple.
\end{remark}

\begin{example}[No real eigenvalues]
Consider $A = \left( \begin{smallmatrix} 0 & -1\\ 1 & 0 \end{smallmatrix} \right)$ over $\R$. Then
\[
\det(A-\lambda\cdot I) = \left| \begin{array}{cc} -\lambda & -1\\ 1 & -\lambda \end{array} \right| = \lambda^2 + 1.
\]
The equation $\lambda^2+1 = 0$ has no solution over $\R$, so $A$ has no real eigenvalues.
\end{example}

\begin{example}[Complex eigenvalues]
The same matrix $A = \left( \begin{smallmatrix} 0 & -1\\ 1 & 0 \end{smallmatrix} \right)$ — the rotation by
$\tfrac\pi2$ — viewed over $\K = \C$ does have eigenvalues: $\lambda^2+1 = 0$ gives $\lambda = \pm i$. For
$\lambda = i$ the system $(A-iI)X=\0$ reads $-ix - y = 0$, so $(1,-i)^T$ is an eigenvector. By the
Fundamental Theorem of Algebra, over $\C$ every square matrix has at least one eigenvalue.
\end{example}

\begin{example}[A multiple eigenvalue]
For the $3\times 3$ identity matrix $I$,
\[
\det(I-\lambda\cdot I) = \left| \begin{array}{ccc} 1-\lambda & 0 & 0\\ 0 & 1-\lambda & 0 \\ 0 & 0 & 1-\lambda \end{array} \right| = (1-\lambda)^3,
\]
so the only eigenvalue is $\lambda = 1$ (with multiplicity $3$).
\end{example}

\begin{fact}
If $\lambda$ is an eigenvalue of $A$, then it is also an eigenvalue of $A^T$.
\end{fact}

\begin{fact}[Trace and determinant]
Let $A$ be an $n\times n$ matrix over $\C$ with eigenvalues $\lambda_1,\ldots,\lambda_n$ listed with
multiplicities. Then
\[
\lambda_1 \lambda_2 \cdots \lambda_n = \det(A), \qquad
\lambda_1+ \lambda_2 +\ldots+ \lambda_n = a_{11}+a_{22}+\ldots+a_{nn}
\]
(the sum of the diagonal entries of $A$ is called the \emph{trace} of $A$). For a $2\times 2$ matrix these
two quantities already determine the characteristic polynomial:
\[
p_A(\lambda) = \lambda^2 - (a_{11}+a_{22})\,\lambda + \det(A).
\]
\end{fact}

\begin{theorem}
Eigenvectors corresponding to pairwise distinct eigenvalues are linearly independent: if
$\lambda_1,\ldots,\lambda_k$ are distinct eigenvalues of $A$ and $\vv_1,\ldots,\vv_k$ are corresponding
eigenvectors, then the set $\{\vv_1,\ldots,\vv_k\}$ is linearly independent.
\end{theorem}

\section{Geometric interpretation}

Eigenvectors of a matrix $A$ are the directions left unchanged (up to scaling) by the linear mapping
$\vv\mapsto A\vv$. The figure below illustrates the action of a shear mapping: the horizontal vector is an
eigenvector (its direction is preserved), whereas the other vector changes direction.

\begin{center}
\includegraphics[scale=0.8]{Eigenvectors.png}
\end{center}

\section{Diagonalisation}

\begin{definition}
An $n\times n$ matrix $A$ over $\K$ is \emph{diagonalisable} if there exists an invertible matrix $P$ such
that $P^{-1}\cdot A\cdot P$ is a diagonal matrix.
\end{definition}

\begin{theorem}
An $n\times n$ matrix $A$ over $\K$ is diagonalisable if and only if $\K^n$ has a basis
$\{\vv_1,\ldots,\vv_n\}$ consisting of eigenvectors of $A$. In that case one may take
$P = (\vv_1\ \vv_2\ \ldots\ \vv_n)$, the matrix with the eigenvectors as columns, and then
\[
P^{-1}\, A\, P = \left( \begin{array}{cccc}
\lambda_1 & 0 & \ldots & 0\\ 0 & \lambda_2 & \ldots & 0\\ \vdots & \vdots & \ddots & \vdots\\ 0 & 0 & \ldots & \lambda_n
\end{array} \right) =: D,
\]
where $\lambda_i$ is the eigenvalue corresponding to the eigenvector $\vv_i$.
\end{theorem}

\begin{remark}
In the language of Chapter~\ref{ch:linear-mappings}: $P$ is the change-of-basis matrix from the
eigenvector basis to the standard basis, and $D = P^{-1}AP$ is the matrix of the operator
$\vv\mapsto A\vv$ written in the eigenvector basis. ``Diagonalisable'' therefore means ``similar to a
diagonal matrix'', and the closing example of the previous chapter, with
$M_S(F) = \left( \begin{smallmatrix} 3 & 0\\ 0 & 1 \end{smallmatrix} \right)$, was secretly a
diagonalisation.
\end{remark}

\begin{corollary}
If an $n\times n$ matrix has $n$ distinct eigenvalues in $\K$, then it is diagonalisable (eigenvectors for
distinct eigenvalues are linearly independent, and $n$ linearly independent vectors form a basis of $\K^n$).
\end{corollary}

\begin{example}[Powers of a matrix]
Diagonalisation makes it easy to compute powers of a matrix: if $A = P D P^{-1}$, then the inner factors
cancel in pairs,
\[
A^k = (PDP^{-1})(PDP^{-1})\cdots (PDP^{-1}) = P\, D^k\, P^{-1},
\]
and the power $D^k$ of a diagonal matrix is computed entrywise on the diagonal. For
$A = \left( \begin{smallmatrix} 2 & 1\\ 1 & 2 \end{smallmatrix} \right)$ we found the eigenvalues $1$ and
$3$ with eigenvectors $(1,-1)^T$ and $(1,1)^T$, so
\[
P = \left( \begin{array}{cc} 1 & 1\\ -1 & 1 \end{array} \right), \qquad
D = \left( \begin{array}{cc} 1 & 0\\ 0 & 3 \end{array} \right), \qquad
P^{-1} = \frac12\left( \begin{array}{cc} 1 & -1\\ 1 & 1 \end{array} \right),
\]
and therefore, for every $k\in\N$,
\[
A^k = P \left( \begin{array}{cc} 1 & 0\\ 0 & 3^k \end{array} \right) P^{-1}
= \frac12 \left( \begin{array}{cc} 3^k+1 & 3^k-1\\ 3^k-1 & 3^k+1 \end{array} \right).
\]
For $k=1$ this indeed recovers $A$.
\end{example}

\begin{warning}
Not every matrix is diagonalisable, even over $\C$: for
$A = \left( \begin{smallmatrix} 1 & 1\\ 0 & 1 \end{smallmatrix} \right)$ the only eigenvalue is
$\lambda = 1$, and its eigenspace $W_1 = \Span{\{(1,0)^T\}}$ is one-dimensional — there is no basis of
$\C^2$ consisting of eigenvectors of $A$.
\end{warning}

\section{Application: PageRank}

Eigenvectors for the eigenvalue $\lambda=1$ underlie Google's original PageRank algorithm for ranking web
pages.

A \emph{web graph} models the web: web pages are nodes and links are directed edges (arrows). For example:
\[
\begin{tikzcd}[row sep=large, column sep=large]
1 \arrow[r] & 2 \\
3 \arrow[u]\arrow[ur] & 4 \arrow[ul]\arrow[l]\arrow[u]
\end{tikzcd}
\]

\begin{definition}
In a web graph $G$, for a vertex $v$ let $l_v$ denote the number of outgoing edges starting at $v$.
\end{definition}

\noindent In the example above $l_1=1,\ l_2=0,\ l_3=2,\ l_4=3$.

\begin{definition}
For a web graph $G$ define its \emph{stochastic matrix} $S$ whose $ij$-entry is
\[
s_{ij} = \frac{1}{l_i} \quad\text{whenever there is a link from $i$ to $j$,} \qquad s_{ij}=0 \text{ otherwise.}
\]
If $l_i=0$, we set $s_{ii}=1$.
\end{definition}

\noindent For the graph above,
\[
S=\left( \begin{array}{cccc}
0 & 1 & 0 & 0\\
0 & 1 & 0 & 0\\
\tfrac{1}{2} & \tfrac{1}{2} & 0 & 0\\
\tfrac{1}{3} & \tfrac{1}{3} & \tfrac{1}{3} & 0
\end{array} \right).
\]

\begin{fact}
For the stochastic matrix $S$ of a web graph $G$:
\begin{itemize}
\item $0\leq s_{ij} \leq 1$ for all $i,j$;
\item $S\cdot (1,\ldots,1)^T = (1,\ldots,1)^T$;
\item $\lambda = 1$ is an eigenvalue of $S$ (and of $S^T$).
\end{itemize}
\end{fact}

Let $w$ be the vector whose $i$-th entry is the probability that a surfer visits web page $i$. Then $w$
satisfies
\[
S^T w = w, \qquad w \text{ has non-negative entries}, \qquad \text{the entries of } w \text{ sum to } 1.
\]

\begin{remark}
In general $S^T$ may have many eigenvectors for $\lambda=1$ satisfying these properties, so the ranking is not
yet well defined. This is fixed by the Google matrix below.
\end{remark}

\begin{definition}[Google matrix]
Let $\alpha$ be the \emph{damping factor} (e.g.\ $\alpha=0.85$) and let $v$ be a \emph{personalization
vector} with non-negative entries summing to $1$. Put
\[
G = \alpha \cdot S + (1-\alpha) \cdot (1,\ldots,1)^T \cdot v^T.
\]
The term $(1-\alpha)\,(1,\ldots,1)^T v^T$ models teleportation (random jumps).
\end{definition}

\begin{fact}
The matrix $G^T$ has a \emph{unique} eigenvector $\pi$ for $\lambda =1$ (that is, $G^T\pi=\pi$) whose entries
are non-negative and sum to $1$. The $i$-th entry of $\pi$ is the \emph{PageRank} of web page $i$.
\end{fact}
